\section{Proof Interface: Barrier, Channel, and Terminal Risk}
\label{sec:criterion}
To prove the risk separation, we verify three geometric tests.
Write $\widetilde T_v$ for the implemented update, including any
allowed error. The following tests apply uniformly over observations
in an event $E$, conditional on the component under consideration.

\begin{theorem}[Barrier, channel, and terminal tests]\label{thm:criterion}
Let $s$ be a unit-direction coordinate and $Z=j$ a component of an
isotropic Gaussian mixture of width $\tau$. Assume
$s(\mu_j)=\xi+d$ with $d>0$ and the following tests hold.
\begin{enumerate}
\item A closed region $W\subseteq\{s\le\xi\}$ contains $x_0$ for
$Z=j$ and satisfies $\widetilde T_v(W)\subseteq W$ for every
restricted variance. This test is required for one component with
$\pi_j>0$.
\item For every component $z$ with $\pi_z>0$, an observation-based allocation reaches a terminal set $G_Z$
within $H$ calls on $E$. A sufficient channel test for a class
needing intermediate calls uses a strip $[s_-,s_E)\times J$
containing its initial state. Each band update either stays in
the strip or crosses into $G_Z$, and
$s^+\ge q_Cs+(1-q_C)s_*$ with $0\le q_C<1$ and $s_*>s_E$.
Then $s_*-(s_*-s_-)q_C^H\ge s_E$ suffices for entry.
For a fixed holding time, the reached part of $G_Z$ must also
be preserved by the remaining band calls.
\item For every such component, the terminal map preserves $G_Z$,
with common bounds $q,\ell,M$ below. For an analysis center $m$
and $x\in G_Z$, it satisfies
$\norm{\widetilde T_{v_L}(x)-m}\le q\norm{x-m}+\ell$ with
$0\le q<1$. The distance to $m$ at switching is at most $M$.
\end{enumerate}
Write $R_m=\E\norm{m-x^\star}^2$, $p=\Prob(E^c)$,
and let $U$ bound $\big(\E\norm{\widehat x-m}^4\big)^{1/2}$ uniformly
in the budget for the policy in tests 2--3. Then
\begin{align*}
 R_{\rm bad}&\ge\pi_j[d\sqrt{\Prob(E\mid Z=j)}-\tau]_+^2,\\
 R_{\rm band}^{1/2}&\le R_m^{1/2}+
 \left[\left(q^n M+\frac{\ell}{1-q}\right)^2+U\sqrt p\right]^{1/2}.
\end{align*}
The first bound holds for any finite budget. The second uses a
budget of $H+n$ with at least $n$ terminal calls on $E$.
\end{theorem}

\begin{proof}
Induction keeps every restricted trajectory in $W$ on $E$ for class
$j$, so its directional distance from $\mu_j$ is at least $d$ there.
The reverse triangle inequality in conditional $L^2$ subtracts at
most $\tau$ for the within-component displacement.
The affine lower comparison forces the strip exit by $H$ calls.
Iterating the terminal estimate bounds the distance to $m$ on $E$.
Cauchy--Schwarz bounds its squared contribution on $E^c$ by
$U\sqrt p$. The $L^2$ triangle inequality gives the upper bound.
\end{proof}

In a sensing eigenbasis, $1-\gamma(B_{ii}+c(v))\ge0$
makes the exact map nondecreasing in its own coordinate, since
$\partial_i\bar\mu_{v,i}=\operatorname{Var}_{w}(\mu_{Z,i})/v$.
For rectangular regions, inward face tests with an error margin
therefore suffice for the set inclusions above, including unbounded
faces. With bounded templates and a positive lower bound on $c(v)$,
the coordinate update gives a budget-uniform moment envelope.
The criterion can be used with other geometries when these tests
are verified. The structured application verifies these tests with explicit margins.
